% Werke X/1 (1917), native156–158, printed149–151; historical remarks by Schlesinger.
\begingroup
\setlength{\headheight}{25pt}
\SourceLeaf{156}\ArithmeticHead{149}{R}{\scriptsize\shortstack{ÄLTESTE UNTERSUCHUNGEN ÜBER\\LEMNISKATISCHE FUNKTIONEN.}}
\NumberedItem{8}
\noindent{\small [S. 54]}
\par\noindent\hspace*{\parindent}Bestimmung der zwei ersten Coefficienten des Ausdrucks
\[\frac1{\surd(1+\sin\varphi^2)}=a+b\cos2\varphi-c\cos4\varphi-d\cos6\varphi\ldots\]
\[\begin{aligned}
a={}&1-\frac12\mathpunct{.}\frac14\mathpunct{.}\frac21
+\frac12\mathpunct{.}\frac34\mathpunct{.}\frac1{16}\mathpunct{.}\frac{4.3}{1.2}\\
&\quad-\frac12\mathpunct{.}\frac34\mathpunct{.}\frac56\mathpunct{.}\frac1{64}\mathpunct{.}\frac{6.5.4}{1.2.3}+\ldots\\
={}&1-\frac{1.1}{2.2}+\frac{1.1.3.3}{2.2.4.4}-\ldots\\
={}&\int\frac{dx}{\surd(1-x^4)}(0\ldots1)\times\frac2\pi=\frac{2A}\pi=0{,}83462
\end{aligned}\]
\[\begin{aligned}
b&=\frac14-\frac14\mathpunct{.}\frac{3.3}{2.6}\mathpunct{.}\frac{5.5}{4.8}+\text{etc.}\\
&=2a-\frac2A.
\end{aligned}\]
\begin{center}\bfseries BEMERKUNGEN.\end{center}
\small
\makeatletter\tagsleft@true\let\veqno\@@leqno\makeatother
\renewcommand{\thefootnote}{*)}
\noindent\hspace*{\parindent}Die Aufzeichnungen [I.] sind mit den Tagebuchnotizen Nr. 50 vom 7. und Nr. 51 vom 8. Januar 1797
in Beziehung zu setzen. Die Artikel [1.], [3.], [4.], [8.] weisen darauf hin, daß \textsc{Gauss} bei seinen Untersuchungen
von dem Studium der \textsc{Euler}schen Abhandlung \textit{De miris proprietatibus curvae elasticae}\footnote{Acta Acad. Petrop. 1782, II (1786), S. 34, Opera omnia, Ser. I, vol. 21, S. 91; die folgenden Zitate beziehen sich auf die Opera.}
ausgegangen ist. Die »curva elastica rectangula« wird durch die Gleichung
\begin{equation}y=\int_0^x\frac{x^2dx}{\sqrt{1-x^4}}\tag{1}\end{equation}
\noindent in rechtwinkligen Koordinaten gegeben; das Integral
\begin{equation}z=\int_0^x\frac{dx}{\sqrt{1-x^4}}\tag{2}\end{equation}
\noindent mißt ihre Bogenlänge. \textsc{Jacob} und \textsc{Johann Bernoulli} hatten gezeigt\renewcommand{\thefootnote}{**)}\footnote{\textsc{Jac. Bernoulli}, Acta erudit. 1694, S. 276 und 336, Opera, Genevae 1744, S. 601 und 608; \textsc{Joh. Bernoulli}, Acta erudit. 1694, S. 394, Opera omnia, Lausannae 1742, I, S. 119.},
daß das Integral (2) auch die Bogenlänge der Lemniskate ausdrückt. \textsc{Gauss} hat in [4.] und ebenso in der Nr. 51 des \textit{Tagebuchs}
das Integral (2) erst nach der curva elastica benannt, dann aber das Wort \textit{elastica} durchgestrichen und \textit{lemniscata} darüber geschrieben.

Die Integrale in den artt. [1.] und [3.] sind zwischen den Grenzen 0 und 1 in bezug auf $p$, zwischen
0 und $\frac\pi2$ in bezug auf $x$ zu nehmen. Setzen wir (vergl. art. [8.] und Werke III, S. 150)
\clearpage
\SourceLeaf{157}\ArithmeticHead{150}{L}{ANALYSIS.\quad NACHLASS.}
\[\int_0^1\frac{dp}{\sqrt{1-p^4}}=A,\qquad\int_0^1\frac{p^2dp}{\sqrt{1-p^4}}=B,\]
\noindent so ist nach \textsc{Euler} (a. a. O. S. 106)
\begin{equation}AB=\frac\pi2.\tag{3}\end{equation}

Der in [1.] angegebene Wert $1{,}311031$ für $A$ steht bei \textsc{Euler}\renewcommand{\thefootnote}{*)}\footnote{Vergl. die Angabe bei \textsc{Gauss}, Werke III, S. 412 (aus Scheda Aa, S. 3).}
a. a. O. S. 104, der in [3.] bei der Berechnung von $2B$ benutzte Wert $B=0{,}599061$ ebenda S. 99.
Wie der von \textsc{Gauss} in [1.] angegebene \textsc{Stirling}sche Wert zeigt, ist der \textsc{Euler}sche Wert von $A$ zu klein,
entsprechend ist der \textsc{Euler}sche Wert von $B$ zu groß\renewcommand{\thefootnote}{**)}\footnote{\textsc{Stirling} gibt a. a. O. S. 57, $B=0{,}59907011\ 736779611$, vergl. \textsc{Gauss}, Werke III, S. 150.}.
Das in [3.] bei dem \textsc{Fuss}schen Werte $1{,}198\pm$ stehende »$=\frac\pi\varpi$« hat \textsc{Gauss} ersichtlich
erst später (mit anderer Schrift und Tinte) hinzugeschrieben; er hat die Bezeichnung
\[\varpi=2\int_0^1\frac{dp}{\sqrt{1-p^4}}\]
\noindent erst später (siehe Werke III, S. 404\renewcommand{\thefootnote}{***)}\footnote{Abdruck eines Zettels (Fh, Kapsel 50), den \textsc{Schering} (Werke III, S. 493) für die älteste Aufzeichnung zur Lehre von den lemniskatischen Funktionen gehalten hat. \textsc{Schering} scheint also weder die \textit{Exercitat. mathem.} noch den \textsc{Leiste} gekannt zu haben. Das letztere geht auch aus der Bemerkung \textsc{Scherings} (Werke III, S. 496, Zeile 19--21) hervor, vergl. die weiter unten im Abschnitt [III.] abgedruckten Aufzeichnungen aus \textsc{Leiste}.}) eingeführt.
Wir machen schon hier auf die große Bedeutung aufmerksam, die dieser Zahl $\frac\pi\varpi$ in der Entwickelung von \textsc{Gauss}’ hierhergehörigen Arbeiten zukommt.
Die Zahl $1{,}2453$ hat \textsc{Gauss} bei [1.] auch ersichtlich später hingeschrieben; wir werden ihre Bedeutung im Abschnitt [II.] kennen lernen.

Die im art. [2.] auftretende Reihe haben wir auch im art. [6] der \textit{Exercitat. mathem.} (S. 141 dieses Bandes)
für die Umkehrung des Integrals (2) gefunden; sie wird hier direkt aus der Reihe
\[x+\frac1{10}x^5+\frac1{24}x^9+\frac5{208}x^{13}+\ldots\]
\noindent für das Integral (2), die z. B. auch in [6.] auftritt, hergeleitet. Zu den Additions- und Multiplikationsformeln
in [4.], der Formel für $s\,4$, d. h. für den sinus lemniscaticus des vierfachen Bogens am Ende von [6.]
so wie für die Reihenentwicklungen in [6.] vergl. man die Werke III, S. 404, 405, 406 abgedruckten Aufzeichnungen\renewcommand{\thefootnote}{$\dagger$)}\footnote{Werke III, S. 405, art. [3.] muß in der Formel für $\cos\operatorname{lemn}2\varphi$ der Nenner lauten $1+2ss-s^4$ statt $1+2ss+s^4$, vergl. oben die Formel für die »Cosehne des doppelten Bogens«.}.

In bezug auf die Bezeichnungen werde folgendes bemerkt. »Sehne eines Bogens« bedeutet die Umkehrung
des Integrals (2), in dem $z$ die Bogenlänge, $x$ den vom Doppelpunkte auslaufenden Fahrstrahl — also eben die Sehne —
für die Lemniskate darstellt. Die Bezeichnung $s$ hat \textsc{Gauss} dann im art. [6.] in sin geändert, später schreibt er dafür
sl oder sin lemn. Die »Sehne des Komplements« ist (vergl. Werke III, S. 404) die zu
\clearpage
\SourceLeaf{158}\ArithmeticHead{151}{R}{\scriptsize\shortstack{ÄLTESTE UNTERSUCHUNGEN ÜBER\\LEMNISKATISCHE FUNKTIONEN.}}
\[\int_0^1\frac{dx}{\sqrt{1-x^4}}-\int_0^x\frac{dx}{\sqrt{1-x^4}}\]
\noindent gehörige Sehne; sie wird mit $c$, cos, später mit cl oder cos lemn. bezeichnet.

Die Formeln für »Sehne der Summe und Differenz« in [4.] stehen bei \textsc{Euler} a. a. O. S. 111. In [7.] ist
$\varphi$ das Integral (2); in der letzten Gleichung dieses Artikels muß die rechte Seite lauten
\[\varphi\frac{-mm+2m+1}{mm+2m-1}.\]

Entwickelungen nach den Kosinus der Vielfachen des Winkels $\varphi$ für Ausdrücke von der Form
$(1-2z\cos\varphi+z^2)^\nu$ oder $(1-\alpha\cos\varphi)^\nu$ sind zuerst von \textsc{Euler}\renewcommand{\thefootnote}{*)}\footnote{\textit{Pièce qui a remporté le prix de l’Académie des Sciences en 1748 sur les inégalités du mouvement de Saturne et de Jupiter} (Paris, 1749).}
in der Störungstheorie angewandt worden. \textsc{Gauss} hat sie aus dem art. 279 des I. Bandes der
\textit{Institutiones calculi integralis}\renewcommand{\thefootnote}{**)}\footnote{L. \textsc{Euleri} Opera omnia, Series I, vol. 11, S. 165 ff.} sicher gekannt.
Die Reihenentwicklungen für die Koeffizienten $a$ und $b$ in [8.] ergeben sich unmittelbar nach dem an der genannten Stelle auseinandergesetzten Verfahren von \textsc{Euler};
die Integraldarstellung von $a$ und $b$ dürfte \textsc{Gauss} erhalten haben, indem er die gefundenen Reihen mit denjenigen verglich, die \textsc{Euler} in der Abhandlung
über die curva elastica für die Integrale $A$ und $B$ gegeben hat\renewcommand{\thefootnote}{***)}\footnote{L. \textsc{Euleri} Opera omnia, Series I, vol. 21, S. 98; bei \textsc{Euler} ist: $A=c$, $B=a$.};
auch der numerische Wert, den \textsc{Gauss} für $\frac{2A}\pi$ gibt, findet sich in dieser Abhandlung \textsc{Eulers} (a. a. O. S. 104).
Daß \textsc{Gauss} schon zu jener Zeit das \textsc{d’Alembert}sche Verfahren der gliedweisen Integration zwischen den Grenzen 0 und $2\pi$ angewandt haben sollte,
um die Integraldarstellung der Koeffizienten zu finden, scheint wenig wahrscheinlich. Vergl. auch die Auszüge aus der Scheda Ac weiter unten
\textit{Zur Theorie des arithmetisch-geometr. Mittels}, Abschnitt [IV.], art. [2.] und Werke III, S. 128.

Als Abfassungszeit der artt. [1.]—[7.] können die ersten Januartage 1797 angesetzt werden; [8.] dürfte wohl etwas später,
aber jedenfalls vor dem 19. März desselben Jahres (vergl. die Bemerkungen zu [II.]) zu datieren sein.
\par\noindent\hfill\textsc{Schlesinger}.\par
\ClosingRule
\clearpage
\endgroup
